% Options for packages loaded elsewhere
\PassOptionsToPackage{unicode}{hyperref}
\PassOptionsToPackage{hyphens}{url}
%
\documentclass[
]{article}
\usepackage{amsmath,amssymb}
\usepackage{lmodern}
\usepackage{iftex}
\ifPDFTeX
  \usepackage[T1]{fontenc}
  \usepackage[utf8]{inputenc}
  \usepackage{textcomp} % provide euro and other symbols
\else % if luatex or xetex
  \usepackage{unicode-math}
  \defaultfontfeatures{Scale=MatchLowercase}
  \defaultfontfeatures[\rmfamily]{Ligatures=TeX,Scale=1}
\fi
% Use upquote if available, for straight quotes in verbatim environments
\IfFileExists{upquote.sty}{\usepackage{upquote}}{}
\IfFileExists{microtype.sty}{% use microtype if available
  \usepackage[]{microtype}
  \UseMicrotypeSet[protrusion]{basicmath} % disable protrusion for tt fonts
}{}
\makeatletter
\@ifundefined{KOMAClassName}{% if non-KOMA class
  \IfFileExists{parskip.sty}{%
    \usepackage{parskip}
  }{% else
    \setlength{\parindent}{0pt}
    \setlength{\parskip}{6pt plus 2pt minus 1pt}}
}{% if KOMA class
  \KOMAoptions{parskip=half}}
\makeatother
\usepackage{xcolor}
\IfFileExists{xurl.sty}{\usepackage{xurl}}{} % add URL line breaks if available
\IfFileExists{bookmark.sty}{\usepackage{bookmark}}{\usepackage{hyperref}}
\hypersetup{
  pdftitle={Effect of SAE on SUA - Gender Mechanisms, Phase 2: Sorting (Issue \#14)},
  hidelinks,
  pdfcreator={LaTeX via pandoc}}
\urlstyle{same} % disable monospaced font for URLs
\usepackage[margin=1in]{geometry}
\usepackage{graphicx}
\makeatletter
\def\maxwidth{\ifdim\Gin@nat@width>\linewidth\linewidth\else\Gin@nat@width\fi}
\def\maxheight{\ifdim\Gin@nat@height>\textheight\textheight\else\Gin@nat@height\fi}
\makeatother
% Scale images if necessary, so that they will not overflow the page
% margins by default, and it is still possible to overwrite the defaults
% using explicit options in \includegraphics[width, height, ...]{}
\setkeys{Gin}{width=\maxwidth,height=\maxheight,keepaspectratio}
% Set default figure placement to htbp
\makeatletter
\def\fps@figure{htbp}
\makeatother
\setlength{\emergencystretch}{3em} % prevent overfull lines
\providecommand{\tightlist}{%
  \setlength{\itemsep}{0pt}\setlength{\parskip}{0pt}}
\setcounter{secnumdepth}{-\maxdimen} % remove section numbering
\usepackage{booktabs}
\ifLuaTeX
  \usepackage{selnolig}  % disable illegal ligatures
\fi

\title{Effect of SAE on SUA - Gender Mechanisms, Phase 2: Sorting (Issue
\#14)}
\author{}
\date{\vspace{-2.5em}}

\begin{document}
\maketitle

Phase 2 of issue \#14: does SAE re-sort girls across schools? Three
tests. (A) Class gender composition: effect of SAE on the share of
female students, with treated schools split by academic vs vocational
type and by pre-reform college-pipeline intensity --- if centralized
choice moves girls toward college-pipeline schools, their share rises
there and falls elsewhere. (B) Commuting: share of students enrolled
outside their commune of residence in 9th grade, by gender --- a
feasible-set test. (C) Composition: average 8th-grade SIMCE of the
class's female and male students --- measured before high school, so it
moves only through who sorts where. Sample, estimator, seed and helpers
identical to \texttt{3\_effect\_sae\_am.Rmd}.

\hypertarget{data-and-treatment}{%
\subsection{Data and treatment}\label{data-and-treatment}}

\hypertarget{headline-sample}{%
\subsection{Headline sample}\label{headline-sample}}

\hypertarget{a.-class-gender-composition-by-school-type-and-pipeline-intensity}{%
\subsection{A. Class gender composition, by school type and pipeline
intensity}\label{a.-class-gender-composition-by-school-type-and-pipeline-intensity}}

One observation per school x cohort over the whole class; the outcome is
the share of female students. Treated schools are split (i) by academic
(HC) vs vocational (TP) type, classified by the majority track of their
pre-2016 enrollment, and (ii) at the within-region median of pre-reform
college-pipeline intensity (mean application share of the school's
2012--2014 classes), taken over treated schools; never-treated private
schools serve as the comparison group in every arm.

\begin{table}[htbp]
   \caption{Effect on Share of Female Students}\label{tab: share female by school type}
   \centerline{\scalebox{0.85}{\begin{tabular}{lcccc}
      \toprule
                   & HC Schools & TP Schools & High Pipeline & Low Pipeline\\
                   & (1) & (2) & (3) & (4)\\
      \midrule
      Treated & 0.002 & 0.007 & 0.000 & 0.008\\
       & (0.007) & (0.007) & (0.004) & (0.010)\\
      \midrule
      Mean (pre-treatment) & 0.540 & 0.474 & 0.558 & 0.496\\
      Observations & 23,128 &  8,397 & 15,730 & 15,411\\
      \bottomrule
   \end{tabular}}}
\end{table}

\hypertarget{b.-enrollment-outside-the-commune-of-residence}{%
\subsection{B. Enrollment outside the commune of
residence}\label{b.-enrollment-outside-the-commune-of-residence}}

Share of the class's female (male) students whose 9th-grade school is in
a different commune than their residence. If centralized choice relaxes
constraints that bind more for daughters, girls' out-of-commune
enrollment rises relative to boys'.

\begin{table}[htbp]
   \caption{Enrollment Outside the Commune of Residence by Gender}\label{tab: gender out of commune}
   \centerline{\scalebox{0.85}{\begin{tabular}{lcc}
      \toprule
      & \multicolumn{2}{c}{Outside Own Commune} \\
      \cmidrule(lr){2-3}
                   & Female & Male \\
                   & (1) & (2)\\
      \midrule
      Treated & -0.012 & -0.008$^{**}$\\
       & (0.008) & (0.004)\\
      \midrule
      Mean (pre-treatment) & 0.262 & 0.255\\
      Observations & 27,708 & 27,240\\
      \bottomrule
   \end{tabular}}}
\end{table}

\hypertarget{c.-baseline-composition-8th-grade-simce}{%
\subsection{C. Baseline composition: 8th-grade
SIMCE}\label{c.-baseline-composition-8th-grade-simce}}

Average 8th-grade SIMCE of the class's female (male) students,
standardized within test year. The score is measured before high school,
so any effect reflects who sorts into the school, not what the school
does --- separating re-sorting from behavior. The 8th-grade test is not
administered every year, so only cohorts with a tested 8th-grade year
contribute.

\begin{table}[htbp]
   \caption{Baseline 8th-Grade SIMCE Composition by Gender}\label{tab: gender simce8 composition}
   \centerline{\scalebox{0.85}{\begin{tabular}{lcccc}
      \toprule
      & \multicolumn{2}{c}{SIMCE-8 Math}
      & \multicolumn{2}{c}{SIMCE-8 Reading} \\
      \cmidrule(lr){2-3}\cmidrule(lr){4-5}
                   & Female & Male & Female & Male \\
                   & (1) & (2) & (3) & (4)\\
      \midrule
      Treated & -0.176 & -0.006 & -0.255$^{*}$ & -0.142\\
       & (0.120) & (0.078) & (0.149) & (0.098)\\
      \midrule
      Mean (pre-treatment) & -0.151 & 0.007 & 0.044 & -0.117\\
      Observations & 27,708 & 27,240 & 27,708 & 27,240\\
      \bottomrule
   \end{tabular}}}
\end{table}

\hypertarget{summary}{%
\subsection{Summary}\label{summary}}

\begin{verbatim}
## -- A. Share female --
\end{verbatim}

\begin{verbatim}
## HC Schools         0.002 (0.007)
## TP Schools         0.007 (0.007)
## High Pipeline      0.000 (0.004)
## Low Pipeline       0.008 (0.010)
\end{verbatim}

\begin{verbatim}
## 
## -- B. Outside own commune --
\end{verbatim}

\begin{verbatim}
## female  -0.012 (0.008) | male  -0.008 (0.004)
\end{verbatim}

\begin{verbatim}
## 
## -- C. SIMCE-8 composition --
\end{verbatim}

\begin{verbatim}
## SIMCE-8 Math     female  -0.176 (0.120) | male  -0.006 (0.078)
## SIMCE-8 Reading  female  -0.255 (0.149) | male  -0.142 (0.098)
\end{verbatim}

\end{document}
